\documentclass{article}
\usepackage[utf8]{inputenc}
\usepackage[T1]{fontenc}
\usepackage{geometry}
\geometry{
    a4paper,
    left=2.5cm,
    right=2.5cm,
    top=2.5cm,
    bottom=2.5cm
}
\usepackage{graphicx}
\usepackage{float}
\usepackage{hyperref}
\usepackage[english]{babel}
\usepackage{amsmath}
\usepackage{amssymb}
\usepackage{array}
\usepackage{booktabs}
\usepackage{caption}
\usepackage{subcaption}
\usepackage{enumitem}
\usepackage{fancyhdr}
\usepackage{titlesec}
\usepackage{url}
\usepackage{textcomp}
\usepackage{xcolor}
\usepackage{listings}
\usepackage{multirow}
\usepackage{longtable}
\usepackage{tabularx}
\usepackage{times}

% Set up page style
\pagestyle{fancy}
\fancyhf{}
\rfoot{\thepage}
\renewcommand{\headrulewidth}{0pt}

% Title formatting
\titleformat{\chapter}[display]{\normalfont\bfseries}{\chaptertitle\thechapter}{0pt}{\Huge}
\titlespacing*{\chapter}{0pt}{50pt}{40pt}

% Document info
\title{Template Document}
\author{Academic Support Agent}
\date{\today}

\begin{document}
\maketitle

%---------------------------------------------------------------
%--------- DADOS B SICOS DO TRABALHO (Preencher todos) ---------
%---------------------------------------------------------------


\titulo{T tulo do Trabalho}    % T tulo do Trabalho em Portugu s

\tituloingles{Title in english}   % T tulo do Trabalho em Ingl s

\palavraschave{1. Palavra 1. 2. Palavra 2. 3. Palavra 3. 4. Palavra 4. 5. Palavra 5.}  % 5 palavras-chave
\keywords{1. Word 1. 2. Word 2. 3. Word 3. 4. Word 4. 5. Word 5}   % 5 key-words

\DEELautor{Nome}{Sobrenome do Autor}                  % Aluno autor do trabalho

\curso{Engenharia El trica}												% Ex: Engenharia Civil

\orientador{Prof. Dr. Orientador}              % Professor Orientador membro da banca 

\membrob{Prof. Dr. Membro da Banca 1}   % Professor membro da banca 
\membroc{Prof. Dr. Membro da Banca 1}   % Professor membro da banca

\local{Paranava }   % Cidade da institui o

%--------------------------------------------------------------------------------------------


\DEELdedicatoria{Dedico este trabalho a todos aqueles que, de alguma forma,\\ auxiliaram para a concretiza o desta etapa. \\ }  % Texto da Dedicat ria (OBS: dedicat ria   opcional)
                                    % Caso n o queira inserir, deixe em branco ( \dedicatoria{} )

\DEELagradecimentos{Agrade o ao IFPR, por fornecer-me as condi es necess rias para conclus o desta importante etapa da minha vida.} 
                                   % Texto dos Agradecimentos (OBS: Agradecimentos   opcional)
                                   % Caso n o queira inserir, deixe em branco ( \agradecimentos{} )
                                  

\DEELepigrafe{\textit{"O dia da crian a   dia da m e,\\ do pai e das professoras,\\ mas tamb m   o dia dos animais.\\ Sempre que voc  olha uma crian a, h  sempre\\ uma figura oculta, que   um cachorro atr s,\\ o que   algo muito importante."} \\ (Dilma Vana Rousseff)} 
                                   % Texto da Ep grafe (OBS: Ep grafe   opcional)
                                   % Caso n o queira inserir, deixe em branco ( \epigrafe{} )


\DEELresumo{Blablablabla yada yada blabla bla} 
                                   % Texto do Resumo (OBS: Resumo   Obrigat rio)

\DEELabstract{Blablablabla yada yada blabla bla} 
                                   % Texto do Abstract (OBS: Abstract   Obrigat rio)



%----- Elementos Pr -Textuais--------------------------
\ElementosPreTextuais
% Insere todos os elementos Pr -Textuais listados (capa, folha de rosto, ...)
												% capa                   % Obrigat rio
												% folhaderosto           % Obrigat rio
												% fichacatalografica     % Obrigat rio
												% folhaaprovacao         % Obrigat rio
												% dedicatoria            % Opcional
												% agradecimentos         % Opcional
												% epigrafe               % Opcional
												% resumo                 % Obrigat rio
												% abstract               % Obrigat rio
\
%--------------------------------------------------------------------------------------------
%\newpage
%---- Outros Elementos Pr -Textuais ----
%--------Lista de figuras--------------------------------------------------------------------


%--------Lista de tabelas--------------------------------------------------------------------



\renewcommand{\cfttablename}{Quadro\space} 
\listof{quadro}{Lista de quadros}


%--------Lista de Siglas e Abreviaturas------------------------------------------------------
% Content from list_siglas would be included here
   % (OBS: Elemento Condicionado   Necessidade) 
%                       % Veja modelo no arquivo list_siglas.tex
%--------Lista de Simbolos e Nota o --------------------------------------------------------
%% Content from list_simbolos would be included here
 % (OBS: Elemento Condicionado   Necessidade)
%                       % Veja modelo no arquivo list_simbolos.tex
%--------Sum rio-----------------------------------------------------------------------------


\thispagestyle{empty}
%--------------------------------------------------------------------------------------------


%----- Elementos Textuais--------------------------
\textual 
% Content from introducao would be included here
				  % Cap tulo 1 - Texto do cap1.tex
%%
% Content from fundamentacao would be included here
				  % Cap tulo 2 - Texto do cap2.tex
%%
% Content from desenvolvimento would be included here
			  % Cap tulo 3 - Texto do cap3.tex
%%
% Content from resultados would be included here
		          % Cap tulo 4 - Texto do cap4.tex
%%
% Content from conclusao would be included here
				      % Cap tulo 5 - Texto do cap5.tex
%%
%% Insira ou retire quantos cap tulos desejar
%%--------------------------------------------------------------------------------------------


%----- Elementos P s-Textuais--------------------------
%
%--------Refer ncias-------------------------------------------------------------------------
\postextual
 % arquivo bibtex com as refer ncias (modelo.bib)

%%--------Apendice-------------------------------------------------
%\apendice                  % Comando que inclui os apendices (sua autoria)
%%
%% Content from apendice_1 would be included here
				  % Apendice 1 - arquivo apendice_1.tex
%%
%% Content from apendice_2 would be included here
         % Apendice 2 - arquivo apendice_2.tex
%%
%% Descomente caso exista a necessidade da inclus o de apendices
%%---------------------------------------------------------------

%%--------Anexos-------------------------------------------------
%\anexo                     % Comando que inclui os anexos (outros autores)
%%
%% Content from anexo_1 would be included here
				  % Anexo 1 - arquivo anexo_1.tex
%%
%% Content from anexo_2 would be included here
          % Anexo 2 - arquivo anexo_2.tex
%%
%% Descomente caso exista a necessidade da inclus o de anexos
%%---------------------------------------------------------------
\end{document}